\section{Related Work}
\label{sec:formatting}

\subsection{AST: Audio Spectrogram Transformer}

The Audio Spectrogram Transformer (AST) that we implemented for interpreting audio data as spectrogram images is a purely attention-based model used for audio classification. It was conceptualized as a way to test whether the reliance on a CNN was necessary in the domain of audio classification, and how domain tasks could be accomplished with a convolution-free model. Historically, CNN-attention hybrids have been used to meet similar goals. It was common to see CNN with self-attention mechanisms layered on top excelling at audio classification and speech \& emotion recognition. The AST mirrors the Vision Transformer (ViT) for vision tasks, using a simpler Transformer architecture with fewer parameters. Notably, while the AST’s input lengths ranged from 1 to 10 seconds, our inputs were consistently 30 seconds. While audio tasks often suffer from a shortage of available data compared to image tasks, AST was able to combat that by using a modified, pretrained Vision Transformer (ViT). This modified version condensed the 3-channel input of ViT into a single channel for AST and adapted positional embeddings for variable-length input, resulting in a much quicker training process when compared to randomized initialization \cite{AST}.

\subsection{Music Recommendation on Spotify using Deep Learning}

When outlining the goals of our project, we used existing music platforms as a baseline for comparing song recommendation methods and results. While Spotify, one of our primary inspirations, has no public documentation on the algorithm it uses for music recommendation, its approach has been speculated to include multiple machine learning algorithms. One of their major algorithm components is Content Filtering, which involves recording and analyzing the habits of existing listeners. Using these learned habits, the platform can more reliably predict the behaviors and listening habits of similar users. This filtering uses two methods: one is an “Exploitative” method that suggests songs based on existing data, preferences, and habits of similar users; another is an “Explorative” method based more heavily on track features which suggests songs that contrast existing preferences. Natural Language Processing is likely another facet of Spotify’s algorithm, wherein text from song titles, album titles, tags, and searches are used to predict user preferences. It is believed that Audio Models are also used to examine the audial properties of songs, incorporating that information into recommendations. A combination of these approaches blend to create an effective song system with high user retention, highlighting both potential parallels and contrasts for our project \cite{spotify}.

\subsection{Free Music Archive}

The Free Music Archive (FMA) dataset that we used for training our model was derived from an open library of the same name. Curated by WFMU, a long-running U.S. radio station, this dataset highlights legally accessible song tracks and has established itself as a large-scale, high-quality library for audio-based projects. It sets itself apart from many other music datasets, such as OMRAS2 and Codaich, by containing the actual audio files. Previously, the standard dataset for Music Information Retrieval (MIR) was GTZAN, which contained only 1000 clips from 10 genres. In contrast, FMA contains over 100,000 clips from over 160 genres. It also contains a wealth of metadata, with metrics from the actual songs in addition to user interactions. It works especially well for Music Genre Retrieval (MGR) due to its feature organization. Features include several sub-genres as well as a built-in genre hierarchy that follows original track annotations. A collection of over 500 pre-computed features also provide valuable insight into track identification. A combination of these attributes make the dataset exceptionally useful for a variety of tasks with the Music domain \cite{FMA}.

\subsection{Learning Transferable Visual Models From Natural Language Supervision}

Contrastive Language-Image Pretraining (CLIP) is a model designed to align image and text embeddings in a purely attention-based format. Its structure inspired the idea to align our audio-based visual representations with other input like lyrics. Its effectiveness can, in part, be attributed to its massive dataset. The development team cultivated a new dataset of 400 million pairs of images and text, forming a diverse collection of information and greatly reducing the risk of over-fitting. Pretraining approaches with both ResNet-50 and a slightly modified version of the Vision Transformer (ViT) were both used for experimentation. The text encoder is a Transformer with architectural modifications. A decayed learning rate, Adam optimizer, and decoupled weight decay regularization were all used during training. Initial hyperparameters were based on a combination of grid \& random searches with manual tuning. The model was ultimately trained to predict text and image pairings, enabling flexible zero-shot transfer that outperformed many comparable approaches like ViR-M and SimCLRv2. The team found success with engineering the text associated with each image to be more specific and context-focused. The model also performed well when it came to representation learning. Though the performance of smaller CLIP models was more reliant on their training dataset, the models scale well and are able to learn a broader spectrum of tasks. A combination of custom architecture and training ultimately made CLIP a robust model that excelled even when compared to human performance \cite{CLIP}.

\subsection{CLAP: Learning Audio Concepts From Natural Language Supervision}

Another model prevalent in the MIR space is the audio alternative to CLIP, CLAP (Contrastive Language-Audio Pretraining). We only discovered CLAP towards the end of the project’s development and were unable to implement it, but it pairs well with our project goals by introducing a clean method of directly aligning audio with text. It aims to emulate the way that humans are able to recognize and contextualize audial input in the same way that CLIP does with images. It has been trained to work with multiple domains, from Sound Event Classification to Speech-related tasks. Their model also had applications in our domain, addressing Music tasks like genres, strokes, and tonics. It employs Zero-Shot predictions, avoiding the limitations of training with class labels and strict predictions at inference time. CLAP utilizes two encoders and contrastive learning in order to bring text and audio data into a multimodal space, essentially combining the manual chaining of our ViT and CLIP embeddings. The two encoders process audio and text, learning the embedding space based on the inputs’ similarities and differences. Once learned, CLAP is capable of flexible class predictions and a variety of downstream tasks \cite{CLAP}.